\section{Introduction} \label{sec:intro}
Large language models (LLMs) have shown considerable promise in assisting chip design~\cite{liu2023chipnemo} and translating natural-language specifications into register-transfer level (RTL) implementations~\cite{zhao2024codev,liu2024rtlcoder,cui2024origen}. Recent LLM-based Verilog-generation systems report functional-correctness rates above 90\% on public benchmarks. VerilogCoder achieves a 94.2\% pass rate on VerilogEval-Human v2~\cite{ho2025verilogcoder}, and MAGE reports 95.7\% pass@1 on VerilogEval v2~\cite{zhao2025mage}. With a 300-node search budget, EvolVE reaches pass rates of 98.1\% on VerilogEval v2 and 92.0\% on RTLLM v2~\cite{hsin2026evolve}. These results reflect complete generation workflows that incorporate tool feedback and iterative search or repair. However, these results do not consistently translate into comparable effectiveness on practical RTL designs. This discrepancy raises concerns about whether current benchmark scores faithfully measure generalization and, consequently, whether they can reliably guide the development of Verilog-generation models.

Data contamination is a critical concern in this setting~\cite{sainz2023nlp,balloccu2024leak}. Publicly available hardware-design data are scarce because production RTL is commonly proprietary and difficult to redistribute. Model training corpora therefore rely heavily on a limited public data ecosystem that includes open-source repositories, educational material, and Internet-accessible design examples. Many RTL benchmarks are assembled from the same ecosystem. VerilogEval, RTLLM, and RealBench provide necessary and valuable evaluation infrastructure~\cite{liu2023verilogeval,lu2024rtllm,jin2025realbench}, but their source data may overlap with corpora collected for model pretraining or fine-tuning. A model can then achieve a high score by reproducing a benchmark solution, or a close derivative, encountered during training rather than by generalizing to an unseen design. The concern is thus not the validity of these benchmarks themselves, but the difficulty of keeping training and evaluation data independent when both draw from a small pool of public RTL.

Reliable detection of such overlap remains difficult. Existing detectors for natural-language data infer membership from token probabilities~\cite{shi2024detecting}, similarity between generated continuations and reference text~\cite{golchin2024time}, or likelihood changes after textual perturbations~\cite{mattern2023membership}. Recurring RTL syntax can make unseen code predictable, potentially confounding regularity with memorization. Code-specific studies show that semantically equivalent variants can challenge text-matching-based detection~\cite{zhang2026syntaxaware}. Similarly, identifier renaming or equivalent RTL restructuring may obscure exposure without changing functionality. RTL benchmarks also pair natural-language specifications with implementations, so exposure to either component warrants consideration. Large-scale evaluations show that membership inference often performs poorly and that distribution shifts between members and non-members can inflate apparent detection success~\cite{duan2024membership}. These findings caution against transferring fixed thresholds across settings without validation. Undisclosed training corpora also limit access to verified membership labels for detector validation and calibration~\cite{shi2024detecting,balloccu2024leak}. Together, these limitations motivate examining conflicting detector estimates for the same model--benchmark pair (Section~\ref{sec:detection}).

In this work, we conduct a systematic study of data contamination in Verilog-generation LLMs. We first evaluate ten representative detectors across multiple models and RTL benchmarks, revealing substantial inconsistencies among their contamination estimates. To establish a reliable basis for comparison, we construct controlled datasets with known membership labels from models whose fine-tuning data are accessible. Confirmed training samples provide contaminated examples, while carefully matched held-out designs provide non-member examples. We then develop a boosting-based detector that learns to integrate complementary signals from the existing detectors. The proposed detector achieves an average AUC of 0.9208 and outperforms every individual detector. Finally, we apply it to representative Verilog-generation LLMs and standard RTL benchmarks. The resulting contamination estimates are approximately 40\%--70\% for most evaluated settings, suggesting that current benchmark scores can substantially overestimate model generalization.

The main contributions of this work are summarized as follows:
\begin{itemize}
\item We systematically evaluate existing contamination detectors in the Verilog domain and demonstrate that their estimates vary substantially across detectors, models, and benchmarks.

\item We construct a controlled evaluation framework with known membership labels, enabling direct and reproducible assessment of contamination detection without assuming that any individual detector is reliable. Our data and code are publicly available on GitHub\footnote{xxx}.

\item We develop a boosting-based detector that integrates complementary contamination signals and achieves an average AUC of 0.9208. Using this detector, we reveal widespread contamination in current Verilog-generation models and demonstrate the need for contamination-aware evaluation protocols.
\end{itemize}
